\subsection*{一步一步建立直观}

\paragraph{1. 把单元看成一条条件化的贡献}

前馈层先把输入投到许多方向，经过非线性，再把各单元贡献合成输出。右图保留一个方向和一个 ReLU：激活强度说明这条方向上的证据有多强。实际向量输出还会乘以一个写出方向，所以一个单元可以理解为“读取一个方向，再写入另一个方向”。



\paragraph{2. 非线性让贡献具有选择性}

当输入位于关闭一侧，当前单元不贡献输出；跨过边界后，贡献随距离线性增长。许多这样的边界可以组合出更复杂的分区。若去掉所有非线性，即便叠加很多仿射层，整体仍只是一个仿射变换，不会自动得到新的表达层次。



\paragraph{3. 关联解释需要结合整体网络}

改变权重和偏置会移动读取方向与门槛，但某个神经元的激活不必对应单一、稳定、易命名的事实。多个特征可能共用方向，一个特征也可能分布在许多单元中。本图用于拆解计算结构；对真实模型的“知识存储位置”仍须做实证检验。
